% ============================================================
% 分节契约
% purpose:      给出"已解决 / 有争议 / 没人做"三分版图,定位本文
% allowed claims: 每条对已有工作的概括须能指回精读卡与 evidence_table
% forbidden claims: 只引不评的流水账;为衬托自己曲解竞品
% inputs:       01-literature/cards/*.md · evidence_table.md
% 验证方式:     nature-citation + nature-ref-verifier 逐条核对
% ============================================================
\section{Related Work}\label{sec:related}

% 按主题分小节组织(不是按时间罗列);每小节末尾一句"与本文的差异"
\subsection{<主题 A(如:消息传递框架)>}
<概述与代表工作 \cite{...}。与本文的差异:……>

\subsection{<主题 B(如:几何/三维方法)>}
<概述与代表工作 \cite{...}。与本文的差异:……>

\subsection{<主题 C(如:数据基准与评估协议)>}
<概述与代表工作 \cite{...}。与本文的差异:……>
